\chapter{Cases and Declension}

\section{The four cases}

German case marks the job performed by a noun phrase. The tables below are a
compact reference; the examples show the same masculine noun changing case.

\begin{center}
\begin{tabularx}{\textwidth}{@{}llllX@{}}
\toprule
Case & Typical function & Question & Masculine example & Example sentence \\
\midrule
\rowcolor{caseNominativeBg}
Nominativ & subject & wer?/was? & der Mann & \textit{Der Mann wartet.} \\
\rowcolor{caseAccusativeBg}
Akkusativ & direct object & wen?/was? & den Mann & \textit{Ich sehe den Mann.} \\
\rowcolor{caseDativeBg}
Dativ & indirect object & wem? & dem Mann & \textit{Ich helfe dem Mann.} \\
\rowcolor{caseGenitiveBg}
Genitiv & possession/association & wessen? & des Mannes & \textit{Das Auto des Mannes.} \\
\bottomrule
\end{tabularx}
\end{center}

\section{Definite and indefinite articles}

\subsection{Definite article}

\begin{center}
\begin{tabular}{@{}lMNFP@{}}
\toprule
 & Maskulin & Neutrum & Feminin & Plural \\
\midrule
\rowcolor{caseNominativeBg}
Nominativ & der & das & die & die \\
\rowcolor{caseAccusativeBg}
Akkusativ & den & das & die & die \\
\rowcolor{caseDativeBg}
Dativ & dem & dem & der & den \\
\rowcolor{caseGenitiveBg}
Genitiv & des & des & der & der \\
\bottomrule
\end{tabular}
\end{center}

\subsection{Indefinite article and \textit{ein}-words}

The same endings are used by \textit{kein} and possessive determiners such as
\textit{mein}, \textit{dein}, \textit{sein}, \textit{ihr}, \textit{unser}, and
\textit{euer}. A dash marks the absence of an indefinite plural article.

\begin{center}
\begin{tabular}{@{}lMNFP@{}}
\toprule
 & Maskulin & Neutrum & Feminin & Plural \\
\midrule
\rowcolor{caseNominativeBg}
Nominativ & ein & ein & eine & -- \\
\rowcolor{caseAccusativeBg}
Akkusativ & einen & ein & eine & -- \\
\rowcolor{caseDativeBg}
Dativ & einem & einem & einer & -- \\
\rowcolor{caseGenitiveBg}
Genitiv & eines & eines & einer & -- \\
\bottomrule
\end{tabular}
\end{center}

\section{Adjective endings}

The determiner and adjective share the job of signalling gender, number, and
case. After a definite article the adjective normally has a weak ending. After
an \textit{ein}-word it supplies information missing from the determiner. With
no determiner, the adjective carries the strong ending itself.

\subsection{Weak endings after the definite article}

\begin{center}
\begin{tabular}{@{}lMNFP@{}}
\toprule
 & Maskulin & Neutrum & Feminin & Plural \\
\midrule
\rowcolor{caseNominativeBg}
Nominativ & -e & -e & -e & -en \\
\rowcolor{caseAccusativeBg}
Akkusativ & -en & -e & -e & -en \\
\rowcolor{caseDativeBg}
Dativ & -en & -en & -en & -en \\
\rowcolor{caseGenitiveBg}
Genitiv & -en & -en & -en & -en \\
\bottomrule
\end{tabular}
\end{center}

\subsection{Mixed endings after an \textit{ein}-word}

\begin{center}
\begin{tabular}{@{}lMNFP@{}}
\toprule
 & Maskulin & Neutrum & Feminin & Plural with possessive/\textit{kein} \\
\midrule
\rowcolor{caseNominativeBg}
Nominativ & -er & -es & -e & -en \\
\rowcolor{caseAccusativeBg}
Akkusativ & -en & -es & -e & -en \\
\rowcolor{caseDativeBg}
Dativ & -en & -en & -en & -en \\
\rowcolor{caseGenitiveBg}
Genitiv & -en & -en & -en & -en \\
\bottomrule
\end{tabular}
\end{center}

\subsection{Strong endings without a determiner}

\begin{center}
\begin{tabular}{@{}lMNFP@{}}
\toprule
 & Maskulin & Neutrum & Feminin & Plural \\
\midrule
\rowcolor{caseNominativeBg}
Nominativ & -er & -es & -e & -e \\
\rowcolor{caseAccusativeBg}
Akkusativ & -en & -es & -e & -e \\
\rowcolor{caseDativeBg}
Dativ & -em & -em & -er & -en \\
\rowcolor{caseGenitiveBg}
Genitiv & -en & -en & -er & -er \\
\bottomrule
\end{tabular}
\end{center}

\subsection{Model noun phrases}

\begin{center}
\begin{tabularx}{\textwidth}{@{}lXXX@{}}
\toprule
 & Definite article & \textit{ein}-word & No article \\
\midrule
\rowcolor{caseNominativeBg}
Nominativ & der gute Wein & ein guter Wein & guter Wein \\
\rowcolor{caseAccusativeBg}
Akkusativ & den guten Wein & einen guten Wein & guten Wein \\
\rowcolor{caseDativeBg}
Dativ & mit dem guten Wein & mit einem guten Wein & mit gutem Wein \\
\rowcolor{caseGenitiveBg}
Genitiv & wegen des guten Weines & wegen eines guten Weines & wegen guten Weines \\
\bottomrule
\end{tabularx}
\end{center}

\section{Personal and reflexive pronouns}

\begin{center}
\begin{tabular}{@{}lcccccc@{}}
\toprule
 & ich & du & er/sie/es & wir & ihr & sie/Sie \\
\midrule
\rowcolor{caseAccusativeBg}
Akkusativ & mich & dich & ihn/sie/es & uns & euch & sie/Sie \\
\rowcolor{caseDativeBg}
Dativ & mir & dir & ihm/ihr/ihm & uns & euch & ihnen/Ihnen \\
\rowcolor{caseAccusativeBg}
Reflexiv Akk. & mich & dich & sich & uns & euch & sich \\
\rowcolor{caseDativeBg}
Reflexiv Dat. & mir & dir & sich & uns & euch & sich \\
\bottomrule
\end{tabular}
\end{center}

\section{Important noun endings}

\begin{itemize}[leftmargin=*]
  \item Dative plural nouns normally add \textbf{-n} when the plural does not
        already end in \textit{-n} or \textit{-s}: \textit{mit den Kindern}.
  \item Masculine and neuter nouns normally add \textbf{-(e)s} in the genitive:
        \textit{des Mannes}, \textit{des Kindes}.
  \item Weak masculine nouns add \textbf{-(e)n} outside the nominative singular:
        \textit{der Student}, but \textit{den/dem/des Studenten}.
\end{itemize}
